% GENERATED FILE. Edit canonical lessons, then run npm run generate:books.
% canonical-source-hash: fnv1a64:f65076561c8192ed
% canonical-lessons: ML-C571-arrakurrappani, ML-C571-cavarrukutta, ML-C571-svikaranamuri, ML-C571-tukkiyituka, ML-C571-nanaykkuka

\chapter{Repairs Around the House}
\label{ch:ml-repairs-around-the-house}
\hlchaptermodality{571}

\begin{chapteropening}
\textbf{By the end of this chapter:} \emph{I can say repairs, a dustbin, a living room, to hang up and to water plants.}

Complete the last lesson of chapter 571: I can say repairs, a dustbin, a living room, to hang up and to water plants.
\end{chapteropening}

\section[\texorpdfstring{\ml{അറ്റകുറ്റപ്പണി} (aṟṟakuṟṟappaṇi)}{അറ്റകുറ്റപ്പണി (aṟṟakuṟṟappaṇi)}]{\ml{അറ്റകുറ്റപ്പണി} (aṟṟakuṟṟappaṇi) — repairs; maintenance}

\label{lesson:ML-C571-arrakurrappani}



\begin{quote}
Before the new one: say the Malayalam for to do the laundry, then the Malayalam for to polish.
\end{quote}

\subsection*{You'll want to know: \ml{അറ്റകുറ്റപ്പണി}}
\textbf{\ml{അറ്റകുറ്റപ്പണി}} — \emph{aṟṟakuṟṟappaṇi} — \textquotedblleft{}repairs; maintenance\textquotedblright{}.

\begin{quote}
\ml{മഴക്കാലത്തിന്} \ml{മുൻപ്} \ml{വീടിന്റെ} \ml{അറ്റകുറ്റപ്പണി} \ml{നടത്തി}.
\end{quote}

\begin{quote}
\emph{maḻakkālattinŭ munpŭ vīṭinṟe aṟṟakuṟṟappaṇi naṭatti.}
\end{quote}

\begin{quote}
Before the rainy season the house was repaired.
\end{quote}

Work on what is broken off and faulty: \ml{അറ്റ} + \ml{കുറ്റ} + \ml{പണി}.

\begin{cousinweb}
\textbf{\ml{അറ്റകുറ്റപ്പണി}} (\emph{aṟṟakuṟṟappaṇi}) is native Dravidian, a compound of \ml{അറ്റ}, \textquotedblleft{}severed\textquotedblright{}, from \ml{അറുക} (DEDR 315), \ml{കുറ്റ}, \textquotedblleft{}faulty\textquotedblright{} (DEDR 1851), and \ml{പണി}, \textquotedblleft{}work\textquotedblright{} (DEDR 3884).
\end{cousinweb}

\begin{grammarlens}[title={What you've built}]
One more word for describing what you lived and felt.
\end{grammarlens}

\subsection*{Your turn}
\begin{itemize}
\raggedright
  \item \emph{Say it:} \emph{aṟṟakuṟṟappaṇi}
  \item \emph{Say it:} \emph{aṟṟakuṟṟappaṇi}, once more
  \item \emph{Say it:} \emph{maḻakkālattinŭ munpŭ vīṭinṟe aṟṟakuṟṟappaṇi naṭatti.} — Before the rainy season the house was repaired.
  \item \emph{Recall:} say the Malayalam for to do the laundry, then the Malayalam for to polish, then say \emph{aṟṟakuṟṟappaṇi} again
\end{itemize}

\begin{tcolorbox}[breakable,colback=teal!4,colframe=teal!35!black,title={Before you move on}]
What does \ml{അറ്റകുറ്റപ്പണി} mean? (\textquotedblleft{}Repairs; maintenance\textquotedblright{}.) Say it once more.
\end{tcolorbox}
\section[\texorpdfstring{\ml{ചവറ്റുകുട്ട} (cavaṟṟukuṭṭa)}{ചവറ്റുകുട്ട (cavaṟṟukuṭṭa)}]{\ml{ചവറ്റുകുട്ട} (cavaṟṟukuṭṭa) — a dustbin}

\label{lesson:ML-C571-cavarrukutta}



\begin{quote}
Before the new one: say the Malayalam for to polish, then the Malayalam for repairs.
\end{quote}

\subsection*{You'll want to know: \ml{ചവറ്റുകുട്ട}}
\textbf{\ml{ചവറ്റുകുട്ട}} — \emph{cavaṟṟukuṭṭa} — \textquotedblleft{}a dustbin\textquotedblright{}.

\begin{quote}
\ml{കുട്ടി} \ml{കടലാസ്} \ml{ചവറ്റുകുട്ടയിൽ} \ml{ഇട്ടു}.
\end{quote}

\begin{quote}
\emph{kuṭṭi kaṭalāsŭ cavaṟṟukuṭṭayil iṭṭu.}
\end{quote}

\begin{quote}
The child put the paper in the dustbin.
\end{quote}

\ml{ചവറ്} rubbish + \ml{കുട്ട} a basket.

\begin{cousinweb}
\textbf{\ml{ചവറ്റുകുട്ട}} (\emph{cavaṟṟukuṭṭa}) is native Dravidian, built from \ml{ചവറ്}, \textquotedblleft{}rubbish, sweepings\textquotedblright{} (DEDR 2395), and \ml{കുട്ട}, \textquotedblleft{}a basket\textquotedblright{}, whose own origin is not settled.
\end{cousinweb}

\begin{grammarlens}[title={What you've built}]
One more word for describing what you lived and felt.
\end{grammarlens}

\subsection*{Your turn}
\begin{itemize}
\raggedright
  \item \emph{Say it:} \emph{cavaṟṟukuṭṭa}
  \item \emph{Say it:} \emph{cavaṟṟukuṭṭa}, once more
  \item \emph{Say it:} \emph{kuṭṭi kaṭalāsŭ cavaṟṟukuṭṭayil iṭṭu.} — The child put the paper in the dustbin.
  \item \emph{Recall:} say the Malayalam for to polish, then the Malayalam for repairs, then say \emph{cavaṟṟukuṭṭa} again
\end{itemize}

\begin{tcolorbox}[breakable,colback=teal!4,colframe=teal!35!black,title={Before you move on}]
What does \ml{ചവറ്റുകുട്ട} mean? (\textquotedblleft{}A dustbin\textquotedblright{}.) Say it once more.
\end{tcolorbox}
\section[\texorpdfstring{\ml{സ്വീകരണമുറി} (svīkaraṇamuṟi)}{സ്വീകരണമുറി (svīkaraṇamuṟi)}]{\ml{സ്വീകരണമുറി} (svīkaraṇamuṟi) — a living room}

\label{lesson:ML-C571-svikaranamuri}



\begin{quote}
Before the new one: say the Malayalam for repairs, then the Malayalam for a dustbin.
\end{quote}

\subsection*{You'll want to know: \ml{സ്വീകരണമുറി}}
\textbf{\ml{സ്വീകരണമുറി}} — \emph{svīkaraṇamuṟi} — \textquotedblleft{}a living room\textquotedblright{}.

\begin{quote}
\ml{അതിഥികൾ} \ml{സ്വീകരണമുറിയിൽ} \ml{ഇരുന്നു}.
\end{quote}

\begin{quote}
\emph{atithikaḷ svīkaraṇamuṟiyil irunnu.}
\end{quote}

\begin{quote}
The guests sat in the living room.
\end{quote}

\ml{സ്വീകരണം} a welcome + \ml{മുറി} a room.

\begin{cousinweb}
\textbf{\ml{സ്വീകരണമുറി}} (\emph{svīkaraṇamuṟi}) is a Malayalam compound built from \ml{സ്വീകരണം}, from Sanskrit \emph{svīkaraṇa}, \textquotedblleft{}making one's own, accepting\textquotedblright{} (Monier-Williams), and \ml{മുറി}, \textquotedblleft{}a room\textquotedblright{}, which is native (DEDR 5008).
\end{cousinweb}

\begin{grammarlens}[title={What you've built}]
One more word for describing what you lived and felt.
\end{grammarlens}

\subsection*{Your turn}
\begin{itemize}
\raggedright
  \item \emph{Say it:} \emph{svīkaraṇamuṟi}
  \item \emph{Say it:} \emph{svīkaraṇamuṟi}, once more
  \item \emph{Say it:} \emph{atithikaḷ svīkaraṇamuṟiyil irunnu.} — The guests sat in the living room.
  \item \emph{Recall:} say the Malayalam for repairs, then the Malayalam for a dustbin, then say \emph{svīkaraṇamuṟi} again
\end{itemize}

\begin{tcolorbox}[breakable,colback=teal!4,colframe=teal!35!black,title={Before you move on}]
What does \ml{സ്വീകരണമുറി} mean? (\textquotedblleft{}A living room\textquotedblright{}.) Say it once more.
\end{tcolorbox}
\section[\texorpdfstring{\ml{തൂക്കിയിടുക} (tūkkiyiṭuka)}{തൂക്കിയിടുക (tūkkiyiṭuka)}]{\ml{തൂക്കിയിടുക} (tūkkiyiṭuka) — to hang up}

\label{lesson:ML-C571-tukkiyituka}



\begin{quote}
Before the new one: say the Malayalam for a dustbin, then the Malayalam for a living room.
\end{quote}

\subsection*{You'll want to know: \ml{തൂക്കിയിടുക}}
\textbf{\ml{തൂക്കിയിടുക}} — \emph{tūkkiyiṭuka} — \textquotedblleft{}to hang up\textquotedblright{}.

\begin{quote}
\ml{അമ്മ} \ml{കഴുകിയ} \ml{തുണികൾ} \ml{അയയിൽ} \ml{തൂക്കിയിട്ടു}.
\end{quote}

\begin{quote}
\emph{amma kaḻukiya tuṇikaḷ ayayil tūkkiyiṭṭu.}
\end{quote}

\begin{quote}
Mother hung the washed clothes on the line.
\end{quote}

Past: \ml{തൂക്കിയിട്ടു}, \emph{tūkkiyiṭṭu}.

\ml{തൂക്കി}, having hung, + \ml{ഇടുക}, to put.

\begin{cousinweb}
\textbf{\ml{തൂക്കിയിടുക}} (\emph{tūkkiyiṭuka}) is native Dravidian, a compound of \ml{തൂക്കുക}, \textquotedblleft{}to suspend, to hang up\textquotedblright{} (DEDR 3376), and \ml{ഇടുക}, \textquotedblleft{}to put\textquotedblright{} (DEDR 442).
\end{cousinweb}

\begin{grammarlens}[title={What you've built}]
One more word for describing what you lived and felt.
\end{grammarlens}

\subsection*{Your turn}
\begin{itemize}
\raggedright
  \item \emph{Say it:} \emph{tūkkiyiṭuka}
  \item \emph{Say it:} \emph{tūkkiyiṭuka}, once more
  \item \emph{Say it:} \emph{amma kaḻukiya tuṇikaḷ ayayil tūkkiyiṭṭu.} — Mother hung the washed clothes on the line.
  \item \emph{Recall:} say the Malayalam for a dustbin, then the Malayalam for a living room, then say \emph{tūkkiyiṭuka} again
\end{itemize}

\begin{tcolorbox}[breakable,colback=teal!4,colframe=teal!35!black,title={Before you move on}]
What does \ml{തൂക്കിയിടുക} mean? (\textquotedblleft{}To hang up\textquotedblright{}.) Say it once more.
\end{tcolorbox}
\section[\texorpdfstring{\ml{നനയ്ക്കുക} (nanaykkuka)}{നനയ്ക്കുക (nanaykkuka)}]{\ml{നനയ്ക്കുക} (nanaykkuka) — to water plants}

\label{lesson:ML-C571-nanaykkuka}



\begin{quote}
Before the new one: say the Malayalam for a living room, then the Malayalam for to hang up.
\end{quote}

\subsection*{You'll want to know: \ml{നനയ്ക്കുക}}
\textbf{\ml{നനയ്ക്കുക}} — \emph{nanaykkuka} — \textquotedblleft{}to water plants\textquotedblright{}.

\begin{quote}
\ml{എല്ലാ} \ml{വൈകുന്നേരവും} \ml{അച്ഛൻ} \ml{ചെടികൾ} \ml{നനയ്ക്കും}.
\end{quote}

\begin{quote}
\emph{ellā vaikunnēravuṁ acchan ceṭikaḷ nanaykkuṁ.}
\end{quote}

\begin{quote}
Every evening Father waters the plants.
\end{quote}

Past: \ml{നനച്ചു}, \emph{nanaccu}.

\ml{നനവുള്ള}, wet, shares the root; it also means to wet: \ml{തുണി} \ml{നനച്ചു}.

\begin{cousinweb}
\textbf{\ml{നനയ്ക്കുക}} (\emph{nanaykkuka}) is native Dravidian: DEDR 3630 gives it as to wet, soak and irrigate (DEDR 3630).
\end{cousinweb}

\begin{grammarlens}[title={What you've built}]
One more word for describing what you lived and felt.
\end{grammarlens}

\subsection*{Your turn}
\begin{itemize}
\raggedright
  \item \emph{Say it:} \emph{nanaykkuka}
  \item \emph{Say it:} \emph{nanaykkuka}, once more
  \item \emph{Say it:} \emph{ellā vaikunnēravuṁ acchan ceṭikaḷ nanaykkuṁ.} — Every evening Father waters the plants.
  \item \emph{Recall:} say the Malayalam for a living room, then the Malayalam for to hang up, then say \emph{nanaykkuka} again
\end{itemize}

\begin{tcolorbox}[breakable,colback=teal!4,colframe=teal!35!black,title={Before you move on}]
What does \ml{നനയ്ക്കുക} mean? (\textquotedblleft{}To water plants\textquotedblright{}.) Say it once more.
\end{tcolorbox}
